
\begin{table}[h]
    \centering
    \begin{tabular}{lr}
        \textbf{Cuenta de Resultados - 5ª Decisión} & \textbf{€} \\
        \hline
        \textbf{Ingresos por ventas}                & \textbf{1183680} \\
        Valor Stock inicial                         & 2297752 \\
        Componentes comprados                       & 38850 \\
        Adquisición de materias primas              & 4055 \\
        Funcionamiento máquinas                     & 73807 \\
        Salarios mecanizado                         & 194687 \\
        Salarios montaje                            & 110812 \\
        Control de calidad                          & 2535 \\
        Transporte                                  & 31400 \\
        Menos valor de stock final                  & 2071121 \\
        \cline{2-2}
        \textbf{Coste de las ventas}                & 682777 \\
        \cline{2-2}
        \textbf{Beneficio Bruto}                    & 500903 \\
        Gastos generales                            & 702089 \\
        Seguros recibidos                           & 0 \\
        Amortización                                & 38454 \\
        \cline{2-2}
        \textbf{EBIT}                               & -239640 \\
        Ingresos financieros                        & 0 \\
        \textbf{Gastos financieros}                 & \textbf{61546} \\
        \cline{2-2}
        \textbf{Beneficio antes de impuestos}       & \textbf{-301186} \\
        Impuestos                                   & 0 \\
        \cline{2-2}
        \textbf{Resultado neto}                     & \textbf{-301186} \\
        \textbf{Beneficio por acción (cents)}       & \textbf{-7,52965} \\
        \textbf{Pago de dividendos}                 & \textbf{0} \\
        \cline{2-2}
        Transferido a reservas                      & -301186 \\
        Reservas acumuladas                         & -2104749 \\
        \cline{2-2}
        \textbf{Reservas totales}                   & \textbf{-2405935}
    \end{tabular}
\end{table}

\begin{table}[h]
    \centering
    \begin{tabular}{|c|r|}
        \hline
        Margen Explotación  & -20,25 \% \\
        Margen Operativo    & -25,44 \% \\
        Margen Neto         & -25,44 \% \\
        ROA                 & -5,11 \% \\
        ROE                 & -18,89 \% \\
        Rotación            & 25,26 \% \\
        Endeudamiento       & 193,99 \% \\
        \hline
    \end{tabular}
\end{table}

\section*{Un Respiro Operativo: Ventas Récord y Vuelta al Margen Positivo}
Tras la asfixia del trimestre anterior, en este tercer trimestre hay un punto de inflexión en la capacidad comercial de IndusMoney. Las medidas correctivas en precios y la inercia del mercado han disparado los ingresos por ventas hasta 1.183.680€, más del doble que en el periodo anterior. Pero la noticia más importante no es solo cuánto se ha vendido, sino cómo.

Por primera vez en lo que va de año, se ha logrado revertir el margen negativo. El Beneficio Bruto ha alcanzado los 500.903€, dejando atrás los números rojos operativos. Esto indica que el ``corazón´´ del negocio ha vuelto a latir: el coste de fabricar los productos ya no supera al precio de venta, lo cual es el primer paso vital para cualquier recuperación, aunque lamentablemente, no ha sido suficiente para frenar los números rojos de la cuenta de resultados.

\section*{Corrección Drástica de Stock: El Grifo Cerrado}
Aprendiendo de los errores del pasado, la dirección de operaciones ejecutó una maniobra de ``frenazo en seco´´ en el aprovisionamiento. La decisión fue clara, ante un sobreabastecimiento de materias primas no se compraron más, dejando el aprovisionamiento en 0 unidades en todos los plazos (este trimestre, el siguiente y su posterior).

Esta decisión fue un acierto necesario para empezar a bajar el gran inventario que se había obtenido en el anterior trimestre. Gracias a esto, y al aumento de las ventas, se ha conseguido reducir ligeramente el inventario de materias primas y, sobre todo, aumentar la rotación de activos al 25\%. Sin embargo, el lastre sigue siendo pesado: todavía quedan más de 2 millones de euros inmovilizados en material, lo que impide que esa liquidez fluya hacia la caja para pagar deudas.

\section*{La Trampa de la Deuda: Del 146\% al 194\% de Endeudamiento}
Si bien la operativa comercial mejoró, la situación financiera ha entrado en una espiral crítica. A pesar de ingresar más de un millón de euros, la caja no ha podido absorber los enormes gastos estructurales y la deuda previa. El resultado es que el descubierto bancario ha crecido hasta los 2.366.244€.

Este aumento descontrolado del pasivo ha provocado que el ratio de endeudamiento suba hasta el 194\%. La empresa debe ahora casi el doble de lo que vale su patrimonio. IndusMoney ha caído en un círculo vicioso: necesita crédito para operar, pero el coste de ese crédito (que ha generado 61.546€ en gastos financieros) consume cualquier beneficio operativo que se genere.

\section*{Los Gastos Generales y la Paradoja de la Producción}
Aunque el margen bruto fue positivo, el Beneficio Neto final sigue siendo negativo con un valor de -301.186€. ¿Dónde se ha ido el dinero? La respuesta está en la estructura de costes fijos y la subcontratación.

Para atender la demanda, se aumentó el nivel de producción a 2 turnos y se contrató a 6 nuevos trabajadores. Sin embargo, simultáneamente se mantuvo una fuerte subcontratación. Esta estrategia fue la misma que en el anterior trimestre, lo que ha disparado los gastos generales hasta los 702.089€, una cifra desorbitada que se ha llevado el beneficio bruto obtenido. Se está produciendo mucho y vendiendo bien, pero la estructura administrativa y operativa es demasiado cara para el tamaño actual de la empresa.

\section*{Patrimonio al Límite y Solvencia Comprometida}
El balance final de la decisión muestra una empresa en estado de alarma financiera. Las pérdidas continuadas han reducido el Patrimonio Neto a 1.594.065€. Aunque los ratios de rentabilidad como el ROE de -19\% y el ROA de -5\% han mejorado ligeramente respecto a la catástrofe del trimestre anterior, siguen en terreno negativo.

El Fondo de Maniobra se mantiene en un nivel alarmante de -1.449.656€. Con una deuda a corto plazo que supera los 2,5 millones de euros y una caja inexistente, IndusMoney afronta la recta final del año dependiendo totalmente de la paciencia de los acreedores. La empresa ha demostrado que sabe vender, pero aún no ha demostrado que sepa ser rentable bajo el peso de su propia deuda.

\section*{La Crónica de una Descoordinación Letal}
La trayectoria de IndusMoney a lo largo de estos cinco trimestres se resume en la incapacidad de compaginar la ambición comercial con la realidad operativa y financiera. La empresa comenzó su declive en las primeras decisiones, se invirtió agresivamente en ``activar la demanda sin asegurar la capacidad´´, generando un escenario de ``publicidad para estanterías vacías´´. El desajuste entre un marketing expansivo y una producción conservadora provocó una ruptura de stock masiva y una caída de ventas del 50\% que erosionó la confianza del mercado y la caja inicial.

Lejos de corregirse suavemente, la gestión sufrió un ``efecto péndulo´´ en el segundo trimestre de 2018. La dirección tomó una decisión errónea, ante la escasez previa realizó un aprovisionamiento masivo y una expansión de fábrica mal financiada. Este error de cálculo transformó una crisis logística en una crisis de solvencia. Se pasó de no tener qué vender a inmovilizar 2,2 millones de euros en stock, activando un descubierto bancario millonario que disparó el endeudamiento por encima del 140\%.

El desenlace en tercer trimestre de 2018 confirma la inviabilidad de IndusMoney. Aunque se logró maquillar la cuenta de resultados con un récord de ventas y la recuperación del margen bruto positivo, el daño estructural ya era irreversible. La empresa ha terminado trabajando exclusivamente para sus acreedores. El peso de una deuda es de 194\%, con unos gastos generales desbordados por la subcontratación. Se cierra el ejercicio demostrando que, en la gestión empresarial, el volumen de ventas es irrelevante si la estructura de capital está rota.

\begin{table}[h]
    \centering
    \begin{tabular}{lr}
        \textbf{Balance - 5ª Decisión}               & \textbf{€} \\
        \hline
        \textbf{Activo fijo}                & \\
        Terreno                             & 50000 \\
        Bienes inmuebles                    & 400000 \\
        Valor de las máquinas               & 1499656 \\
        \cline{2-2}
        \textbf{Total activo fijo}          & \textbf{1949656} \\
        \textbf{Activo Circulante}          & \\
        Inventario - Productos              & 7106 \\
        Inventario - Componentes            & 38850 \\
        Inventario - Materia Prima          & 2025165 \\
        Deudores                            & 665560 \\
        Efectivo e Inversores               & 0 \\
        \cline{2-2}
        \textbf{Total activo circulante}    & \textbf{2736681} \\
        \cline{2-2}
        \textbf{Activo Total}               & \textbf{4686337} \\
        \textbf{Pasivo:}                    & \\
        Impuestos a pagar                   & 0 \\
        Acreedores                          & 226028 \\
        Línea de crédito / Descubierto      & 2366244 \\
        \cline{2-2}
        \textbf{Pasivo exigible}            & \textbf{2592272} \\
        \textbf{Préstamos}                  & \textbf{500000} \\
        \textbf{Activo neto}                & \textbf{1594065} \\
        \textbf{Fondos propios}             & \\
        Capital social                      & 4000000 \\
        Prima de emisión                    & 0 \\
        Reservas                            & -2405935 \\
        \cline{2-2}
        \textbf{Patrimonio neto}            & 1594065 
    \end{tabular}
\end{table}

\begin{table}[h]
    \centering
    \begin{tabular}{|c|r|}
        \hline
        NOF & 144.409,00 € \\
        FM & -1.449.656,00 € \\
        EC & -1.594.065,00 € \\
        \hline
    \end{tabular}
\end{table}